% AY2018 制御工学 解答例
% 回答の流れ・言葉遣いは 03_制御工学/参考/「制御工学Ⅱ（完成版）」に寄せる。
% デザイン（囲み枠等）は真似しない。解答・数値は本年度過去問から自力で導いた。
\documentclass[11pt,a4paper]{article}
\usepackage{../../../00_common/preamble}

\usetikzlibrary{positioning,arrows.meta,calc}
\usepackage{pgfplots}
\pgfplotsset{compat=newest}

\title{制御工学 AY2018 解答例（専門応用科目I）}
\author{}
\date{}

\begin{document}
\maketitle

% ==================================================================
\section*{第1問〔I〕}

質量 $m$ にバネ $k$ とダンパ $d$ を並列に介して入力変位 $x_1$ を与える系。$x_2$ は質量の位置。

\subsection*{(1) 運動方程式と伝達関数}
バネ・ダンパは質量と入力端 $x_1$ を結ぶので、質量に働く力は $k(x_1-x_2)$ と $d(\dot x_1-\dot x_2)$。
\begin{align*}
  &\qquad m\ddot{x}_2=k(x_1-x_2)+d(\dot{x}_1-\dot{x}_2)\\
  &\therefore\quad m\ddot{x}_2+d\dot{x}_2+kx_2=d\dot{x}_1+kx_1\\
  &\therefore\quad (ms^2+ds+k)X_2=(ds+k)X_1.
\end{align*}
\[
  \therefore\quad \ans{\;G(s)=\dfrac{X_2(s)}{X_1(s)}=\dfrac{ds+k}{ms^2+ds+k}\;}
\]

\subsection*{(2) ステップ応答（$m=1,d=3,k=2$）}
$G(s)=\dfrac{3s+2}{s^2+3s+2}=\dfrac{3s+2}{(s+1)(s+2)}$。$X_1(s)=\dfrac1s$ より
\begin{align*}
  &\qquad X_2(s)=\frac{3s+2}{s(s+1)(s+2)}
   =\frac{1}{s}+\frac{1}{s+1}-\frac{2}{s+2}\\
  &\therefore\quad x_2(t)=1+\e^{-t}-2\e^{-2t}.
\end{align*}
直流ゲイン $G(0)=\dfrac{2}{2}=1$ なので定常値は $1$。
\[
  \therefore\quad \ans{\;x_2(t)=1+\e^{-t}-2\e^{-2t},\qquad x_2(\infty)=1\;}
\]

\subsection*{(3) 最大値と発生時間}
\begin{align*}
  &\qquad \dot{x}_2(t)=-\e^{-t}+4\e^{-2t}=0
   \quad\therefore\quad 4\e^{-2t}=\e^{-t}\\
  &\therefore\quad \e^{t}=4
   \quad\therefore\quad t=\ln 4.
\end{align*}
このとき $x_2=1+\tfrac14-2\cdot\tfrac{1}{16}=\tfrac98$。分子の零点によりオーバーシュートを生じる。
\[
  \therefore\quad \ans{\;x_{2\max}=\dfrac98\ \ (t=\ln4\approx1.39)\;}
\]

\subsection*{(4) $k=0$ の場合（$m=1,d=1,k=0$）}
\begin{align*}
  &\qquad G(s)=\frac{s}{s^2+s}=\frac{s}{s(s+1)}=\frac{1}{s+1}\\
  &\therefore\quad X_2(s)=\frac{1}{s(s+1)}=\frac1s-\frac{1}{s+1}
   \quad\therefore\quad x_2(t)=1-\e^{-t}.
\end{align*}
$G(0)=1$ より定常値は $1$。バネがないため定常状態では質量が入力変位に追従する。
\[
  \therefore\quad \ans{\;x_2(t)=1-\e^{-t},\qquad x_2(\infty)=1\;}
\]

\bigskip
\noindent 検算: (2) は $x_2(0)=0$・定常 $1$、(3) は $t=\ln4$ で $x_{2\max}=\tfrac98$、
(4) は $G=1/(s+1)$ で定常 $1$ を SymPy で確認✓。

\newpage
% ==================================================================
\section*{第2問〔II〕}

$\frac1M\to\frac1s\to\frac1s\to\alpha$ の直列に、速度帰還 $D$（第1積分器の後段）と
位置帰還 $K$（第2積分器の後段）をもつフィードバック系。

\subsection*{(1) 伝達関数}
加速度を $A$、速度を $V=\frac1s A$、変位 $P=\frac{A}{s^2}$、出力 $Y=\alpha P$ とすると
\begin{align*}
  &\qquad A=\frac1M\Bigl[U-KP-DV\Bigr]
   =\frac1M\Bigl[U-\frac{K}{s^2}A-\frac{D}{s}A\Bigr]\\
  &\therefore\quad A\Bigl(M+\frac{D}{s}+\frac{K}{s^2}\Bigr)=U
   \quad\therefore\quad Y=\frac{\alpha}{s^2}A=\frac{\alpha\,U}{Ms^2+Ds+K}.
\end{align*}
\[
  \therefore\quad \ans{\;G(s)=\dfrac{\alpha}{Ms^2+Ds+K}\;}
\]

\subsection*{(2) ゲインと位相（$M=5,D=5.5,K=0.5,\alpha=5$）}
\begin{align*}
  &\qquad G(s)=\frac{5}{5s^2+5.5s+0.5}=\frac{10}{(10s+1)(s+1)}.
\end{align*}
$s=j\omega$ として、直流ゲイン $|G(0)|=10$（$20\,\mathrm{dB}$）。
\[
  \therefore\quad \ans{\;|G(j\omega)|=\dfrac{10}{\sqrt{1+(10\omega)^2}\,\sqrt{1+\omega^2}},\quad
  \angle G(j\omega)=-\tan^{-1}10\omega-\tan^{-1}\omega\;}
\]

\subsection*{(3) ゲイン線図（漸近線）}
$|G(0)|=20\,\mathrm{dB}$。折れ点は $10s+1$ より $\omega=0.1$、$s+1$ より $\omega=1\,[\mathrm{rad/s}]$。
傾きは $\omega<0.1$ で $0$、$0.1<\omega<1$ で $-20$、$\omega>1$ で $-40\,[\mathrm{dB/dec}]$。

\begin{center}
\begin{tikzpicture}
\begin{axis}[
  width=0.86\linewidth, height=5.2cm,
  xmin=1e-2, xmax=1e2, ymin=-80, ymax=40, xmode=log,
  xlabel={$\omega\,[\mathrm{rad/s}]$}, ylabel={ゲイン [dB]},
  ytick={-80,-60,-40,-20,0,20,40},
  xmajorgrids=true, ymajorgrids=true, xminorgrids=true,
  grid style={line width=0.3pt, draw=gray!60},
  extra y ticks={0}, extra y tick style={grid=major, grid style={line width=1pt, draw=black}},
  legend pos=south west, clip=true, enlargelimits=false,
]
\addplot[thick] coordinates {(0.01,20) (0.1,20) (1,0) (100,-80)};
\addlegendentry{漸近線}
\addplot[thick,dashed] coordinates
  {(0.01,20.0) (0.1,17.0) (0.316,7.0) (1,-3.0) (10,-40.1) (100,-80)};
\addlegendentry{真値}
\node[anchor=south] at (axis cs:0.1,20) {\scriptsize $\omega=0.1$};
\node[anchor=north east] at (axis cs:1,0) {\scriptsize $\omega=1$};
\end{axis}
\end{tikzpicture}
\end{center}

\subsection*{(4) 位相進み補償後の極と安定判別}
$G_c(s)=\dfrac{1+0.25s}{1+0.025s}$ を直列結合すると
\begin{align*}
  &\qquad G(s)G_c(s)=\frac{10}{(10s+1)(s+1)}\cdot\frac{1+0.25s}{1+0.025s}.
\end{align*}
極は分母 $(10s+1)(s+1)(1+0.025s)=0$ の根である。
\begin{align*}
  &\qquad 10s+1=0,\quad s+1=0,\quad 1+0.025s=0\\
  &\therefore\quad s=-0.1,\ -1,\ -40.
\end{align*}
\[
  \therefore\quad \ans{\;\text{極}\ s=-0.1,\,-1,\,-40\ \text{はすべて実部が負}\;\Rightarrow\;\text{安定}\;}
\]
補償の零点 $s=-4$ は位相を進めるが極は上記のとおりで、3 極とも左半平面にあるので
閉ループは安定である。

\bigskip
\noindent 検算: (2) は $G=10/((10s+1)(s+1))$、$|G(0)|=10$。(4) は極 $\{-0.1,-1,-40\}$ が
すべて Re$<0$ で安定であることを SymPy で確認✓。

\end{document}
